\documentclass[11pt]{article}
\usepackage{fontspec}
\usepackage[french,english]{babel}
\usepackage{amsmath}
\usepackage{amssymb}
\usepackage{amsthm}
\usepackage{mathrsfs}
\usepackage{tikz-cd}
\usepackage{enumitem}

\newcommand{\printedpage}[1]{\marginpar{\footnotesize p.~#1}}
\newcommand{\msp}[1]{\marginpar{\footnotesize [#1]}}
\newcommand{\Tr}{\mathrm{Tr}}
\newcommand{\Cr}{\mathrm{Cr}}
\newcommand{\Hom}{\mathrm{Hom}}
\newcommand{\Ty}{\underline{\mathrm{Ty}}}
\newcommand{\Ob}{\mathrm{Ob}}
\newcommand{\lng}{\operatorname{long}}
\newcommand{\card}{\operatorname{card}}
\newcommand{\bD}{\boldsymbol{\Delta}}
\newcommand{\uTy}{\underline{\mathrm{Ty}}}

\begin{document}

\printedpage{24}
% head: A. Grothendieck : Dérivateurs, Ch. VII, Édition des Universités de Montpellier II et Paris VII

\textbf{Corollary 3.} \textit{If $f$ is an} increasing \textit{morphism of oriented intervals, then for every $J' \in \Cr(\tau')$, we have}
\[
(0 \leq)\ \card \Cr(f)^{-1}(\{J\}) \longrightarrow \lng(J'),
\] % [sic] \longrightarrow -> \leq, the next sentence speaks of equality; and {J} -> {J'}, as in Cr(f)^{-1}({J'}) below
\msp{page 37}
\textit{equality being obtained only if the parts of $f$ on the critical intervals of $I_\tau$ which lie over $J'$ are all equal to $0$ or $1$ (hence if and only if on the $J \in \Cr(f)^{-1}(\{J'\})$, they are equal to $1$).}

\textbf{Corollary 4.} \textit{Suppose moreover that $I'$} \texttt{[}$= I_{\tau'}$\texttt{]} \textit{is reduced. Then}
\[
\mathit{Cr}(f) : \mathit{Cr}(\tau, f) \longrightarrow \mathit{Cr}(\tau')\ (\simeq \mathrm{Fl}\,\mathrm{tr}(\tau'))
\]
\textit{is} bijective.

Thus in this case, $\tau'$ reduced, we obtain a map in the opposite direction
\[
\begin{tikzcd}[column sep=small]
\Tr(\tau') \simeq \Cr(\tau') \arrow[rr, "\Cr^*(f)"] \arrow[dr, "\Cr(f)^{-1}"'] & & \Cr(\tau) \\
 & \Cr(\tau, f), \arrow[ur, hook, "\text{inclusion}"'] &
\end{tikzcd}
\]
a map which transforms the critical interval $J'$ $(= \{i - 1, i\})$ into the \textit{unique} critical interval $J$ of $I_\tau$ such that
\[
f(J) = J'.
\]

\msp{page 38}
But $J \simeq \bD_i^{\varepsilon}$, and we therefore have a \textit{surjective} morphism
\[
f_0 : \bD_i^{\varepsilon} \longrightarrow J' \simeq \bD_1^{\varepsilon}.
\]
We see right away that there then exists a \textit{unique} transition arrow $u$ of $\bD_i^{\varepsilon}$, such that $f(u) \neq \mathrm{id}$. Thus we find the following corollary, more precise than corollary 4:

\textbf{Corollary 5.} \textit{Let $f : I_\tau \longrightarrow I_{\tau'}$ be an increasing morphism of oriented intervals, with $I_{\tau'}$ reduced, and let $\Tr(\tau, f)$ be the set of transition arrows $\underline{u}_i$ of $I_\tau$ such that $f(\underline{u}_i) \neq \mathrm{id}$, i.e. such that $\pi(u_i, f) = 1$. Then $f$ induces a map} % [sic] telle -> telles, agreement with "flèches" (cf. "telles que" just after)
\[
\Tr(\tau, f) \longrightarrow \Tr(\tau')
\]
\textit{which is bijective, and (of course) increasing, hence an isomorphism of ordered sets.}

Thus the map $\Tr(\tau') \longrightarrow \Cr(\tau, f)$ from earlier is made more precise by a bijective map
\[
\Tr(f)^* : \Tr(\tau') \longrightarrow \Tr(\tau, f),
\]
related to $\Cr(f)^*$ above by the commutative triangle
\[
\begin{tikzcd}
 & \Tr(\tau, f) \arrow[dd, "\wr\ \alpha"] \\
\Tr(\tau') \arrow[ur, "\Tr(f)^*"] \arrow[dr, "\Cr(f)^*"'] & \\
 & \Cr(\tau, f),
\end{tikzcd}
\]
% foot: Version 22 décembre 2001
\printedpage{25}
% head: A. Grothendieck : Dérivateurs, Ch. VII, Édition des Universités de Montpellier II et Paris VII
\msp{page 39}
the map $\alpha$ associating to every $u_i \in \Tr(\tau, f)$ the unique critical interval which ``contains'' $u_i$ (i.e. such that $u_i$ is an arrow of the said interval, or again, is part of the associated ``critical chain'' \dots).

We have better still:

\textbf{Proposition 8.} \textit{Let $\tau, \tau' \in \mathrm{Ty}$, $\tau'$} reduced. \textit{For every increasing morphism of oriented intervals}
\[
f : I_\tau \longrightarrow I_{\tau'}
\]
\textit{(we shall simply say that $f$ is an arrow in $\uTy$, and write $f \in \Hom_{\uTy}(I_\tau, I'_\tau)$), let us associate to it the subset $\Tr(\tau, f)$ of transition arrows of $I_\tau$ nondegenerate for $f$. These arrows form a totally ordered set equipped with signs $+$ or $-$ (direct or retrograde character of the arrows), and $f$ induces a bijection of totally ordered sets} % [sic] I'_\tau -> I_{\tau'}, the target of f is I_{\tau'}
\[
\underbrace{E_f}_{\subseteq \Tr(\tau)} = \Tr(\tau, f) \simeq \Tr(\tau'),
\]
\textit{and even an isomorphism of totally ordered sets ``equipped with signs''. Thus, to every $f$ we associate a subset $E_f$ of $\Tr(\tau)$, and such that its ``type'' is equal to $\tau'$. Let us denote by}
\[
\mathfrak{P}_{\tau'}(\Tr(\tau)) \subseteq \mathfrak{P}(\Tr(\tau))
\]
\textit{the set of subsets $E$ of $\Tr(\tau)$ whose type is $\tau'$.}

\msp{page 40}
\textit{Then the map}
\[
f \longmapsto \Tr(\tau, f) : \Hom_{\uTy}(\tau, \tau') \longrightarrow \mathfrak{P}_{\tau'}(\Tr(\tau))
\]
\textit{is} bijective \textit{: giving oneself a morphism $f : \tau \longrightarrow \tau'$ in $\uTy$,} amounts to the same thing \textit{as giving oneself a strictly increasing sequence of transition arrows}
\[
\underline{u}_{i_1},\ \underline{u}_{i_2},\ \dots\ ,\underline{u}_{i_{n'}} \qquad (n' = \lng \tau')
\]
\textit{in $I_\tau$ $(i_1 < i_2 < \cdots < i_{n'})$, such that the corresponding sequence of signatures}
\[
(\varepsilon_{i_1}, \varepsilon_{i_2}, \dots, \varepsilon_{i_{n'}}) = \tau'
\]
\textit{is equal to $\tau'$.}

\textbf{Corollary 1.} \textit{For every $\tau$, we have canonical bijections}
\[
\left\{
\begin{array}{lcl}
\Tr^+(\tau) & \simeq & \Hom_{\uTy}(\tau, \Delta^+) \\
\Tr^-(\tau) & \simeq & \Hom_{\uTy}(\tau, \Delta^-),
\end{array}
\right.
\]
functorial in $\tau$ \textit{(varying in} $\uTy$\textit{).}
% foot: Version 22 décembre 2001

\printedpage{26}
% head: A. Grothendieck : Dérivateurs, Ch. VII, Édition des Universités de Montpellier II et Paris VII

To give the corollary a meaning, we have to define
\[
\tau \longmapsto \Tr^+(\tau), \qquad \tau \longmapsto \Tr^-(\tau)
\]
as \textit{functors} contravariant in $\tau$, thus for a morphism
\[
f : I_{\tau_2} \longrightarrow I_{\tau_1}
\]
in $\uTy$, define maps
\msp{page 41}
\[
\begin{array}{rcl}
\Tr^+(f)^* : \Tr^+(\tau_1) & \longrightarrow & \Tr^+(\tau_2) \\
\Tr^-(f)^* : \Tr^-(\tau_1) & \longrightarrow & \Tr^-(\tau_2),
\end{array}
\]
in other words, a map
\[
\Tr(f)^* : \Tr(\tau_1) \quad \longrightarrow \quad \Tr(\tau_2)
\]
``compatible with the signatures''. In corollary 5 above, we defined such a map only for $\tau'$ reduced \texttt{[}$\tau_1 = \tau'$\texttt{,} $\tau_2 = \tau$\texttt{]} (\footnote[12]{What follows does not involve the signatures - it is a sorites on increasing maps between finite totally ordered sets.}). For an arrow $u$ of $I_\tau$, let us call \textit{set-theoretic part} of $f$ on $u$, where of $u$ for $f$, the set-theoretic part of the arrow $f(u)$ in $I_{\tau'}$, i.e. the segment $[f(i-1), f(i)] \subseteq I_{\tau'}$ (here $u = \underline{u}_i$). Let % [sic] où -> ou, alternative "de f sur u, ou de u pour f"
\[
\underline{u}_{i_1},\ \dots\ ,\underline{u}_{i_N} \qquad \text{or } (v_1, v_2, \dots, v_n)
\] % [sic] v_n -> v_N, the sequence has N elements (cf. v_N below)
be the increasing sequence of elements of $\Tr(\tau, f)$, we then see by induction on $i$ that for $i \in [1, N-1]$, we have
\[
\begin{array}{rcl}
f(\mathrm{but}(v_i)) & = & f(\mathrm{source}(v_{i+1}))\ \ (= t_i), \\
f(\mathrm{source}(v_1)) & = & 0_{\tau'}, \\
f(\mathrm{but}(v_N)) & = & 1_{\tau'}.
\end{array}
\]
Let us consider the sequence
\[
\{\overbrace{0_{\tau'}}^{= t_0}, t_1, t_2, \dots, t_{N-1}, \overbrace{1_{\tau'}}^{= t_N}\},
\]
it is defined by the corresponding subset of $I_{\tau'}$, which we shall call the set of \textit{critical vertices of $I_{\tau'}$ for $f$.} It is a sequence containing $0_{\tau'}$, $1_{\tau'}$, and such that two consecutive elements are connected in $I_{\tau'}$ by an arrow, which is equivalent to saying that it contains the critical vertices (intrinsic) of $I_{\tau'}$. On the other hand, for every transition interval $[i-1, i]$ of $I_{\tau'}$ $(i \in [1, n'])$, % [sic] appelera -> appellera, French spelling
% foot: Version 22 décembre 2001
\printedpage{27}
% head: A. Grothendieck : Dérivateurs, Ch. VII, Édition des Universités de Montpellier II et Paris VII
\msp{page 42}
there exists one and only one transition interval $[t_{j-1}, t_j]$ for this sequence $(j \in [1, N])$, such that we have $[i-1, i] = [t_{j-1}, t_j]$. This says that there exists one and only one $v_j \in \Tr(\tau, f)$ such that the part of $f$ on $v_j$ (namely precisely $[j-1, j]$) contains (that of) $\underline{u}_i$. To sum up:

\textbf{Corollary 2.} \textit{Let $f : \tau \longrightarrow \tau'$ be in $\uTy$. Then for every $u \in \Tr(\tau') = \Tr(I_{\tau'})$, there exists one and only one $v \in \Tr(\tau)$ such that the part of $f$ on $v$ contains $u$. Moreover, the signatures of $u$ and $v$ are equal.}

It is in this way that we therefore obtain a map
\[
\Tr(f)^* : \Tr(\tau') \quad \longrightarrow \quad \Tr(\tau) \qquad (\text{factoring through } \Tr(\tau, f))
\]
inducing
\[
\begin{array}{rcl}
\Tr^+(f)^* : \Tr^+(\tau') & \longrightarrow & \Tr^+(\tau) \\
\Tr^-(f)^* : \Tr^-(\tau') & \longrightarrow & \Tr^-(\tau).
\end{array}
\]
Transitivity for a composite
\[
\tau \xrightarrow{f} \tau' \xrightarrow{g} \tau''
\]
or
\[
I_\tau \xrightarrow{f} I_{\tau'} \xrightarrow{g} I_{\tau''}
\]
is immediate. If $u$, $u'$, $u''$ are transition arrows in $I_\tau$, $I_{\tau'}$, $I_{\tau''}$ respectively, such that
\[
\begin{array}{rcll}
\pi(g, u') & \supseteq & \pi(u''), & \qquad \text{i.e. } u' = \Tr(g)^*(u'') \\
\pi(f, u) & \supseteq & \pi(u'), & \qquad \text{i.e. } u = \Tr(f)^*(u'),
\end{array}
\]
then we have
\[
\pi(gf, u) \ \supseteq\ \pi(u''), \qquad \text{i.e. } u = \Tr(gf)^*(u'').
\]

\msp{page 43}
If the arrows in question have as intrinsic part $[i-1, i]$, $[i'-1, i']$, $[i''-1, i'']$ respectively, the hypothesis is
\[
\begin{array}{rclrcl}
g(i'-1) & \leq & i''-1, & g(i') & \geq & i'' \\
f(i-1) & \leq & i'-1, & f(i) & \geq & i'
\end{array}
\]
and since $g$ is increasing, we have deduced
\[
\begin{array}{rcll}
gf(i-1) & \leq & g(i'-1) & (\leq i''-1) \\
gf(i) & \geq & g(i') & (\geq i''),
\end{array}
\]
\begin{flushright}
q.e.d.
\end{flushright}
% foot: Version 22 décembre 2001

\printedpage{28}
% head: A. Grothendieck : Dérivateurs, Ch. VII, Édition des Universités de Montpellier II et Paris VII

Thus, corollary 1 above acquires a precise meaning and it is proved at the same time, since the correspondences of corollary 1 are given by
\[
\underbrace{f}_{\in \Hom_{\uTy}(\tau, \Delta^{\varepsilon})} \qquad \longmapsto \Tr(f)^*(\underline{u}),
\]
where $\underline{u}$ is the unique transition arrow of $\Delta^{\varepsilon}$.

We are going to use these sorites to give, for every $\tau \in \uTy$, a simple combinatorial description of the categories $\uTy/\tau$ and $\tau\backslash\uTy$, or at least of their sets of objects.

\textbf{1$^o$ Description of} $\Ob\,(\tau\backslash\uTy)$\textbf{.}

To $f \in \Ob(\tau\backslash\uTy)$, i.e.
\[
f : \tau \longrightarrow \tau' \qquad \text{i.e. } f : I_\tau \longrightarrow I_{\tau'}
\]
in $\uTy$, is associated $E = E_f = \Tr(\tau, f) \subseteq \Tr(\tau)$, the set of critical transition arrows (i.e. nondegenerate) for $f$. To each $u \in E$, let us associate the integer $\lng(\pi(f, u))$,
\msp{page 44}
the \textit{part} of $f$ on $u$. Then $(\tau', f)$ is uniquely determined by the pair
\[
(E, \pi), \qquad E \in \mathfrak{P}(\Tr(\tau)), \qquad \underbrace{\pi}_{\text{weighting}} : E \ \longrightarrow\ \underbrace{\mathbf{N}^*}_{\substack{= \text{set of} \\ \text{integers} \geq 1}} \quad .
\]
Or if one prefers, by the map $\pi = \pi_f$
\[
\Tr(\tau) \xrightarrow{\pi} \mathbf{N}, \qquad u \longmapsto \pi(u) = \pi(u, f) = \lng(\pi_{\tau'}(f(u)))
\]
(which gives back $E$ as
\[
E = \pi^{-1}(\mathbf{N}^*) \qquad ).
\]
This map is subject to no restriction,
\[
\Ob\,(\tau\backslash\uTy) \ \xrightarrow[\sim]{\text{bij.}}\ \Hom_{\mathrm{Ens}}(\Tr(\tau), \mathbf{N}).
\]
We recover $\tau'$ from $\pi$ as follows. We take $E = \pi^{-1}(\mathbf{N}^*)$ (transition arrows nondegenerate for $\pi$), $E$ is totally ordered by the order induced by that of $\Tr(\tau)$, it is moreover signatured \texttt{[sic]} by the signaturation \texttt{[sic]} of $\Tr(\tau)$, let
\[
(\varepsilon_1, n_1),\ \dots\ ,\ (\varepsilon_r, n_r)
\]
be the sequence of signatures and weights, of the sequence of elements of $E$ $(\varepsilon_i \in \{\pm 1\}$, $n_i \in \mathbf{N}^*$ for $i \in [1, r])$. Then,
\[
\tau' = \Delta_{\varepsilon_1}^{n_1} \Delta_{\varepsilon_2}^{n_2} \cdots \Delta_{\varepsilon_r}^{n_r}, \qquad r = \card E,
\]
in the free monoid $\mathrm{Ty}$, where we have denoted by $\Delta_+$, $\Delta_-$ (instead of $\Delta^+$, $\Delta^-$) the two generators. The map
\[
I_\tau \ \longrightarrow\ I_{\tau'}
\]
% foot: Version 22 décembre 2001

\end{document}
